\documentclass[11pt]{article}
\usepackage[margin=1in]{geometry}
\usepackage{amsmath,amssymb,bm}
\title{Charge on conductor surfaces and component interfaces}
\author{}
\date{}
\begin{document}
\maketitle

\section{Problem Setup}
The conduction and exterior calculations give an electric field
$\bm E$ on every triangle. We want to know where charge must be to
make that field. The app assumes the same permittivity $\epsilon$
in every material and in air. Within a uniform conductor region,
\begin{equation}
 \nabla\cdot\bm J=0,\qquad
 \bm J=\gamma(\bm E+\bm F).
\end{equation}
Both $\gamma$ and the battery driving field $\bm F$ are constant within
that region, so $\nabla\cdot\bm E=0$ there. Air also has
$\nabla\cdot\bm E=0$. Under this model, charge is represented on
conductor--air surfaces and material interfaces, rather than as
a density within the regions.

\section{Charge from the jump in electric field}
Consider a very thin box crossing an interface. Let its two broad
faces have area $A$, and let unit normal $\bm n$ point from side 1
to side 2. Gauss's law for the electric displacement
$\bm D=\epsilon\bm E$ says that the outgoing flux equals the
enclosed charge:
\begin{equation}
 A\bm D_2\cdot\bm n-A\bm D_1\cdot\bm n=\sigma A.
\end{equation}
As the box thickness shrinks, its side flux vanishes. Dividing by
$A$ gives the surface charge density
\begin{equation}
 \sigma=(\bm D_2-\bm D_1)\cdot\bm n
       =\epsilon(\bm E_2-\bm E_1)\cdot\bm n.
 \label{eq:jump}
\end{equation}
This equation applies both where a conductor meets air and where
two conducting materials meet. The cavity-facing surface is also
a conductor--air interface. The distant rectangle is only the
edge of the computational domain, not a charged surface.

\section{Evaluating the charge on mesh edges}
The potential is piecewise linear, so the electric field is
constant on each triangle, as explained in \texttt{fem.tex}.
For each tagged physical edge $e$, the mesh supplies the two
neighboring triangles $h_1$ and $h_2$ and a normal $\bm n_e$
from the first toward the second. The code evaluates
\begin{equation}
 \sigma_e
 =\epsilon(\bm E_{h_2}-\bm E_{h_1})\cdot\bm n_e.
\end{equation}
This gives one constant density for each tagged edge. The
corresponding charge per unit depth on that edge is
$q'_e=\sigma_e\ell_e$, where $\ell_e$ is its length.
No second field solve is needed. Jumps between triangles of
the same material are not labeled as physical charge sheets.

\paragraph{Code:} \texttt{calculations/charges.py}.
\end{document}
